\section{Threats to Validity}
\label{sec:threats}



The main threats to the validity of our work consist of internal, construct, and external threats.

\textbf{Threats to internal validity} include the influence of the method used to identify the code elements (e.g., function calls, assignments, or control structures). To reduce this threat, we use Joern~\cite{joern} code analyzer, which is widely used in existing studies~\cite{velvet,vuldeeppeaker,ivdetect,linevd}. 

\textbf{Threats to construct validity} relate to the suitability of our evaluation procedure. We used \textit{exact match ratio}, \textit{hamming score}, \textit{accuracy}, \textit{macro-average}, \textit{micro-average}, \textit{weighted-average}, and \textit{sample-average}. They are the classical evaluation measures for multi-label classification~\cite{multilabel_cls}. 
%
% Another threat to the construct validity is 
% The dataset might not be representative.

\textbf{Threats to external validity} mainly lie in the selection of multi-label classification models used in our experiments. To mitigate this threat, we select the representative models which are well-known for NLP and SE tasks. \base is considered as the baseline approach for general multi-label text classification, while all Word2vec, Glove, and CodeBERT are reputed and shown to be effective in many SE tasks. The dataset used in our experiments might not be representative or not very high-quality. To reduce this threat, we used the largest public dataset~\cite{bigvul}, which is collected from a large number of real-world projects and widely used in existing vulnerability detection studies~\cite{linevd, linevul, ivdetect}. Additionally, our data has only C/C++ code. Thus, we cannot claim that similar results would have been observed in other programming languages. Further studies are needed to validate and generalize our findings to other languages. 

% A threat relates to our selections for baseline approaches. To reduce this threat, we selected the state-of-the-art $n$-gram with nested cache, CodeT5~\cite{codet5}, and the combination of GPT-2~\cite{gpt2} and a type-based post-processing step. Moreover, $n$-gram with nested cache is used to represent for the classical direction which is widely used in existing studies~\cite{bigcode,allamanis2018survey,failures}. Meanwhile, CodeT5 and GPT-2 are used as a representative advanced DNN approach instead of vanilla LSTM or RNN. Moreover, GPT-2 is used as the core engine of Tabnine~\cite{tabnine} and IntelliCode~\cite{intellicode}, which are two of the most popular AI-assisted code completers.